% ---------------------------------------------------------------- title page
\pagenumbering{roman}
\begin{titlepage}
\thispagestyle{empty}
\begin{center}
{\color{cetpgrey}\large PES University \,$\cdot$\, Capstone Project (UE23CS441A)\par}
\vspace{2.2cm}
{\fontsize{54}{60}\selectfont\bfseries\color{cetpblue} CETP\par}
\vspace{0.4cm}
{\Large\bfseries Cross-Environment Runtime Prediction\\ for PostgreSQL Queries\par}
\vspace{0.8cm}
{\color{cetpteal}\rule{0.55\textwidth}{1.2pt}\par}
\vspace{0.8cm}
{\LARGE Master Technical Reference\par}
\vspace{0.3cm}
{\large Architecture, implementation, evidence, limitations and defense preparation\par}
\vspace{2.4cm}
\begin{tabular}{@{}r@{\hspace{1.2em}}l@{}}
\rowcolors{1}{white}{white}
\textbf{Document version} & 1.0 (first complete edition)\\
\textbf{Prepared} & 27 October 2026\\
\textbf{Repository state} & branch \texttt{x86-probe}, HEAD \texttt{38c4bf2}, working tree modified\\
\textbf{Software version} & \texttt{cetp} 0.4.0 (classifier ``Alpha'')\\
\textbf{Project team} & Harshith J, Ishanvee Amit Sinha, Sarthak Wadhwa, Kartik Bhat Sumbly\\
\textbf{Guide} & Prof.\ Nivedita Kasturi\\
\textbf{Status} & Research prototype; evidence-based, not a marketing document\\
\end{tabular}
\end{center}
\vfill
{\small\color{cetpgrey}Every quantitative claim in this document carries a source path and an evidence tag
(see the Evidence methodology section above). Where the repository contradicts itself, the contradiction is reported
rather than resolved silently.}
\end{titlepage}

% ---------------------------------------------------------------- status and scope
\chapter*{Project status and document scope}
\addcontentsline{toc}{chapter}{Project status and document scope}
\markboth{Project status and document scope}{}

\textbf{Project status.} CETP is an academic capstone prototype. It predicts how long a PostgreSQL
query will take on a \emph{production} machine from an \texttt{EXPLAIN (ANALYZE)} captured on a
\emph{development} machine plus a description of each machine's capability. The repository contains
two generations of the method (\cref{sec:two-gens}); the installable package is version 0.4.0 and is
classified ``Alpha'' in \src{pyproject.toml}. The repository's own README states: ``The research is
sound; treat the product as alpha.'' This document agrees with the first half of that sentence only
in a restricted sense, spelled out in \cref{ch:exec,ch:limits}.

\textbf{What this document covers.} The state of the repository on 27 October 2026 (HEAD
\texttt{38c4bf2}, branch \texttt{x86-probe}, 17 commits, working tree modified), including the
uncommitted fix of evidence-integrity finding C10 described in \cref{sec:c10}.

\textbf{What it does not cover.} (i) Any measurement taken after that date. (ii) Content of the ca.\ 50
PDF reports under \src{docs/} (reorganised into \texttt{docs/\{design,guides,reference,results\}} in the uncommitted working tree). They were not read; in the audit pass I extracted text from seven of them (CANONICAL, operator\_model, laptop\_to\_prod, cross\_isa\_prediction, HANDBOOK, LLD, DATASHEET) only to look up specific numbers, and equivalence with their deleted Markdown sources is unverified. (iii) The C
sources of the Generation-1 hardware probe (only the header comment of \src{hw\_probe\_suite.c} was read) and the bodies of about 25 analysis scripts, which were not
read in full. (iv) The original EC2 raw logs, which no longer exist in the repository. Each gap is
tracked in \cref{app:issues}.

\tableofcontents
\listoffigures
\listoftables

% ---------------------------------------------------------------- abbreviations
\chapter*{List of abbreviations and terminology}
\addcontentsline{toc}{chapter}{List of abbreviations and terminology}
\markboth{Abbreviations}{}
\small
\begin{longtable}{@{}L{2.6cm}L{12.3cm}@{}}
\toprule \thead{Term} & \thead{Meaning in this document}\\ \midrule \endhead
AURC & Area under the risk--coverage curve; average error when the most confident predictions are kept first.\\
CI & Confidence interval; here always a 95\% percentile cluster-bootstrap interval unless stated.\\
Cluster bootstrap & Resampling whole cross-validation folds (machines) with replacement, so the 420 pairs are not treated as independent.\\
CPU quota (CFS) & Linux Completely Fair Scheduler limit on CPU time (\texttt{docker --cpus}); slows parallel work but not a single thread.\\
CV (timing) & Coefficient of variation of repeated timings, standard deviation divided by mean. Not cross-validation; that is always spelled out.\\
DOP & Degree of parallelism available to a query: $1+\texttt{max\_parallel\_workers\_per\_gather}$ (3 at PostgreSQL defaults).\\
EBS & Amazon Elastic Block Store; network-attached storage used by all five EC2 machines.\\
EXPLAIN ANALYZE & PostgreSQL command that executes a query and returns its plan with actual per-node timings, row counts and buffer counts.\\
Gather & PostgreSQL plan node that launches parallel workers; every node below it runs once per participant.\\
GBR & Gradient-boosted regression trees (scikit-learn \texttt{GradientBoostingRegressor}).\\
GHA & GitHub Actions; its x86\_64 runners supplied the cross-ISA target machines.\\
Gen 1 / Gen 2 & The roofline-style two-scalar method and the per-operator method, respectively (\cref{sec:two-gens}).\\
ISA & Instruction-set architecture; here arm64 (Apple M1) versus x86\_64.\\
LOMO, LOQO, LOMQO & Leave-one-machine-out, leave-one-query-out and leave-one-machine-and-query-out cross-validation.\\
MAPE, MdAPE & Mean and median absolute percentage error (\cref{sec:metrics}).\\
Naive & The do-nothing baseline: predict that production runtime equals development runtime (scaling factor $=1$).\\
Naive-linear & Baseline scaling by $1/\text{bandwidth ratio}$.\\
PMU & Hardware performance-monitoring unit; blocked on AWS Nitro instances, which motivated the DB-only scope.\\
pp & Percentage points.\\
$q$-error & $\max(\hat y/y,\,y/\hat y)$; symmetric multiplicative error used in the database literature.\\
Roofline & Performance model bounding throughput by compute or memory bandwidth; the project's ``roofline'' is a weighted blend, not a true minimum of bounds.\\
Scaling factor & $s = T_{\text{prod}}/T_{\text{dev}}$; the quantity every predictor in this project estimates.\\
SF1 & TPC-H scale factor 1 (about 1\,GB of base data).\\
SLA & Service-level agreement; here a per-query latency budget in milliseconds.\\
STREAM & Memory-bandwidth microbenchmark (triad kernel) \citep{mccalpin:stream}.\\
TPC-H & Decision-support benchmark of 22 analytical queries \citep{tpch:spec}.\\
Warm-up repetition & The first of 20 repetitions of a measurement cell; discarded because it is a cold-cache run.\\
\bottomrule
\end{longtable}
\normalsize

% ---------------------------------------------------------------- purpose and audience
\pagenumbering{arabic}
\setcounter{page}{1}
\chapter*{Document purpose and intended audience}
\addcontentsline{toc}{chapter}{Document purpose and intended audience}
\markboth{Purpose and audience}{}

This is a technical master reference, written so that one person can read it and then:
explain the project to a beginner, a professor, an interviewer or a teammate; trace an input through the
code to the printed prediction; defend each design decision with evidence; reproduce the headline
numbers; and answer detailed follow-up questions, including the uncomfortable ones. It is deliberately
not a brochure. The repository is unusually candid about its own negative results, and this document
keeps that candour.

\begin{center}\small
\begin{tabularx}{\textwidth}{@{}L{3.1cm}Y@{}}
\toprule \thead{Reader} & \thead{Start here}\\ \midrule
Beginner or examiner new to the topic & \cref{ch:exec,ch:problem,ch:background}, then \cref{ch:conclusion}\\
Technical reviewer / professor & \cref{ch:arch,ch:impl,ch:algo,ch:results,ch:limits}\\
Teammate taking over the code & \cref{ch:repo,ch:impl,ch:setup,ch:testing}, \cref{app:reference}\\
Interviewer or viva panel & \cref{ch:exec,ch:results,ch:limits}, and the question bank in \cref{app:defense}\\
Anyone quoting a number & \cref{app:evidence} first, then \src{results/CANONICAL.md}\\
\bottomrule
\end{tabularx}
\end{center}

% ---------------------------------------------------------------- evidence methodology
\chapter*{Evidence and verification methodology}
\addcontentsline{toc}{chapter}{Evidence and verification methodology}
\markboth{Evidence methodology}{}
\label{sec:evidence}

\textbf{Authority order.} Current code and verified execution output are authoritative for how the
system behaves \emph{today}. Result files (\src{results/*.json}, \src{results/*.csv}) are authoritative
for what was measured. Prose in comments, README, PDFs or earlier chats is used only to explain intent
and is never allowed to override the first two. Where prose and measurement disagree, the
disagreement is reported (\cref{app:issues}).

\textbf{Evidence tags.} Each claim in a table or numbered list carries one of these tags.
\begin{center}\small
\begin{tabularx}{\textwidth}{@{}L{1.6cm}Y@{}}
\toprule \thead{Tag} & \thead{Meaning}\\ \midrule
\evRUN & I executed it or recomputed it myself in this session (output seen).\\
\evPACK & Executed by the author of the audited evidence pack (\src{docs/research/}) and recorded there; not re-run by me.\\
\evFILE & Read from a committed result, data or configuration file; the computation behind it was not re-run.\\
\evCODE & Established by reading the source (path and line numbers given).\\
\evCLAIM & Stated in a comment, README or report but not traced to code or data.\\
\evUNK & Unknown or unavailable; the missing step is named.\\
\bottomrule
\end{tabularx}
\end{center}

\textbf{What was actually done for this edition.}
\begin{enumerate}
\item Read in full: all twelve modules of \src{cetp/} (about 2,100 lines), \src{gate/cetp_gate.py},
  \src{scripts/collection/operator_probe.py}, the parts of \src{scripts/analysis/canonical_evaluation.py}
  and \src{operator_model_evaluation.py} that define data loading, predictors, metrics, cross-validation and
  bootstrap, \src{scripts/analysis/build_calibration.py}, the Makefile, \src{pyproject.toml},
  \src{requirements.txt}, \src{Dockerfile}, the README and the three CI workflows (headers).
\item Read the ten evidence-pack files \src{docs/research/00--09*.md} and the committed result files cited in
  \cref{app:evidence}.
\item \evRUN\ \texttt{python3 -m cetp.cli selftest} printed \texttt{selftest PASSED} (5 checks) in a
  Python 3.10.12 Linux environment with no third-party packages.
\item \evRUN\ \texttt{pytest tests/test\_operator\_model.py tests/test\_product\_loop.py}: \textbf{80 passed}
  in 0.10\,s (same environment). The other two suites need scikit-learn, numpy and pandas and were not run by me.
\item \evRUN\ Decomposed and priced the three real EXPLAIN fixtures (\texttt{q1}, \texttt{q18}, \texttt{q19}) with the
  shipped code and the committed operator registry (\cref{sec:worked}).
\item \evRUN\ Recomputed pooled conformal coverage from \src{results/tier1_conformal.csv} (\cref{sec:conformal}).
\item Searched the web to confirm that six of the cited database papers exist with the stated venues.
\end{enumerate}

\textbf{What was not done.} No PostgreSQL server, Docker daemon or EC2 instance was used; no probe or
collection script was run; \texttt{make results} and \texttt{make verify-determinism} were not run by me
(the evidence-pack author ran \texttt{canonical\_evaluation.py} and found its outputs byte-identical to the
committed ones, \evPACK). The full 119-test suite result is \evPACK. The Generation-1 C probe source and the
Chrome-based PDF build were not exercised. The bibliography was written from knowledge of the literature;
only the six papers named above were looked up.

\begin{honest}
Every limitation above is an \emph{unverified} item, not a failed one. Nothing in this document claims that
an unrun command succeeds. Commands in \cref{ch:setup} are labelled \evRUN, \evPACK\ or ``inferred''.
\end{honest}
